\section{Introduction}
\label{sec:introduction}

The global apparel industry generates substantial fabric waste,
a significant portion of which arises from inefficient placement of 2D garment
panels (commonly called \emph{pattern pieces}) on the fabric roll prior
to cutting.
The task of arranging these panels to minimise wasted material is known as
\emph{nesting} (or \emph{marker making}), and has been extensively studied as
a combinatorial optimisation problem~\cite{BennellO08}.
Classical nesting methods, as well as the prevailing industry standard,
treat texture as irrelevant: panels are placed purely for compactness,
with no regard for whether the fabric pattern (stripes, checks, herringbone,
or other periodic motifs) continues seamlessly across the seams of the
assembled garment.

The problem of aligning periodic textures along garment seams has been
addressed by Wolff et al.~\cite{WolffS19}, who directly deform
the 2D patch shapes to achieve alignment.
Their formulation, however, assumes an unbounded fabric supply: roll width and
material efficiency play no role, so the method is decoupled from the
practical nesting problem.
Furthermore, by operating directly on 2D patch geometry, fit changes induced
by the deformation are implicit and go unmeasured. There is no mechanism
to bound or quantify how much the modified pattern deviates from the
originally designed garment.

We argue that seam placement is best optimised in 3D, on the body surface,
rather than by deforming flat patches.
Moving a seam line on the body mesh changes the 2D pattern only as a
consequence of flattening, keeping the modification grounded in 3D geometry.
Critically, this makes fit deviation explicit: the difference between
the flattened and the original 3D surface can be directly measured and
minimised as part of the optimisation.
We formalise this observation and build on it to define a joint optimisation
problem that simultaneously addresses fabric efficiency, seam-aware texture
alignment, and fit preservation, objectives that have not previously been
considered together.

\paragraph{Contributions.}
\begin{enumerate}
    \item We formulate garment pattern layout as a joint optimisation over
          seam positions, fabric efficiency, and texture alignment under
          realistic fabric roll constraints, the first method to address
          all three objectives simultaneously.

    \item We propose a 3D-grounded seam optimisation paradigm in which fit
          deviation is an explicit, measurable objective. In contrast to prior
          2D patch deformation approaches, fit change is no longer implicit
          and unmeasured.

    \item We support texture alignment across a set of wallpaper symmetry
          classes, including patterns with glide-reflection symmetry such as
          herringbone, chevron, and brick bond.

    \item We demonstrate scalability to layouts of up to TODO pattern pieces
          and robustness across TODO garment types (shirts, pants, onesies,
          and dresses).
\end{enumerate}

\paragraph{Outline.}
Section~\ref{sec:related_work} reviews related work on nesting, texture
alignment, and evolutionary approaches to garment optimisation.
Section~\ref{sec:method} presents the problem formulation and the full
optimisation pipeline.
Section~\ref{sec:experiments} reports experimental results.
Section~\ref{sec:conclusion} concludes.
